% Problem set 1: uniformly accelerated straight-line motion (15 July 2026)
% One source, two documents:
%   for son: pdflatex -jobname=2026-07-15-kinematics-uniform-acceleration-for-son kinematics-01-uniform-acceleration.tex
%   for dad: pdflatex -jobname=2026-07-15-kinematics-uniform-acceleration-for-dad "\def\isdad{1}% Problem set 1: uniformly accelerated straight-line motion (15 July 2026)
% One source, two documents:
%   for son: pdflatex -jobname=2026-07-15-kinematics-uniform-acceleration-for-son kinematics-01-uniform-acceleration.tex
%   for dad: pdflatex -jobname=2026-07-15-kinematics-uniform-acceleration-for-dad "\def\isdad{1}% Problem set 1: uniformly accelerated straight-line motion (15 July 2026)
% One source, two documents:
%   for son: pdflatex -jobname=2026-07-15-kinematics-uniform-acceleration-for-son kinematics-01-uniform-acceleration.tex
%   for dad: pdflatex -jobname=2026-07-15-kinematics-uniform-acceleration-for-dad "\def\isdad{1}% Problem set 1: uniformly accelerated straight-line motion (15 July 2026)
% One source, two documents:
%   for son: pdflatex -jobname=2026-07-15-kinematics-uniform-acceleration-for-son kinematics-01-uniform-acceleration.tex
%   for dad: pdflatex -jobname=2026-07-15-kinematics-uniform-acceleration-for-dad "\def\isdad{1}\input{kinematics-01-uniform-acceleration.tex}"
% The son's copy omits the [given -> find] tags and the missing-variable analysis:
% identifying the combination is part of the training.
\documentclass[11pt,a4paper]{article}
\usepackage[margin=2.2cm,bottom=2.4cm]{geometry}
\usepackage[T1]{fontenc}
\usepackage[utf8]{inputenc}
\usepackage{lmodern}
\usepackage{amsmath}
\usepackage{xcolor}
\usepackage{booktabs}
\usepackage{enumitem}
\usepackage{parskip}
\usepackage{fancyhdr}
\usepackage{tcolorbox}
\usepackage{tikz}
\usepackage{url}

\definecolor{accent}{HTML}{2F6FED}
\definecolor{muted}{HTML}{5B6473}
\definecolor{ink}{HTML}{1F2430}
\definecolor{boxbg}{HTML}{F2F6FF}
\definecolor{rule}{HTML}{E2E5EA}

\ifdefined\isdad
  \newcommand{\copyname}{Dad's copy: tags, answers, references}
  \newcommand{\footname}{Dad's copy}
  \newcommand{\given}[1]{\hfill{\small\color{muted}[#1]}}
\else
  \newcommand{\copyname}{student's copy}
  \newcommand{\footname}{student's copy}
  \newcommand{\given}[1]{}
\fi

\pagestyle{fancy}
\fancyhf{}
\renewcommand{\headrulewidth}{0pt}
\renewcommand{\footrulewidth}{0.4pt}
\fancyfoot[L]{\small\color{muted} Uniformly accelerated motion --- problem set 1 (\footname)}
\fancyfoot[R]{\small\color{muted} Page \thepage}

\newcommand{\nsec}[1]{\vspace{9pt}\par{\color{accent}\large\bfseries #1}\par\vspace{2pt}}
\newcommand{\ms}{\ \mathrm{m/s}}
\newcommand{\mss}{\ \mathrm{m/s^2}}

\begin{document}\color{ink}

{\huge\bfseries Uniformly Accelerated Motion}\\[4pt]
{\color{muted} Kinematics of straight-line motion with constant acceleration
\;\textperiodcentered\; problem set 1 \;\textperiodcentered\; \copyname}\\[6pt]
{\color{rule}\rule{\linewidth}{0.9pt}}

\begin{tcolorbox}[colback=boxbg, colframe=accent, boxrule=0.9pt, arc=2.5mm,
                  left=10pt, right=10pt, top=7pt, bottom=7pt,
                  title={\large\bfseries The rules of the game}, coltitle=white]
For \textbf{every} problem, including the warm-up:
\begin{enumerate}[label=\textcolor{accent}{\textbf{\arabic*.}}, leftmargin=1.8em, itemsep=1.5pt, topsep=3pt]
  \item Choose the positive direction and set $x_0 = 0$ at the start.
  \item Solve for what is asked.
  \item Draw the three graphs against time: \textbf{$a(t)$, $v(t)$, $x(t)$} --- axes labelled,
        key values marked (initial/final speeds, times, distances).
\end{enumerate}
Graphs need not be to exact scale, but the \emph{shapes} must be right: with constant
acceleration $a(t)$ is a horizontal line, $v(t)$ is a straight line with slope $a$,
and $x(t)$ is a parabola (bending up if $a>0$, down if $a<0$). Use $g = 10\mss$.
\smallskip

\ifdefined\isdad
\textbf{The toolbox} --- four equations, each \emph{missing} one quantity:
\begin{center}
\renewcommand{\arraystretch}{1.5}
\begin{tabular}{@{}llll@{}}
\toprule
& \textbf{Equation} & \textbf{Missing} & \textbf{Use when\ldots}\\
\midrule
(i)   & $v = u + at$                    & $s$ & distance does not matter\\
(ii)  & $s = ut + \tfrac{1}{2}at^{2}$   & $v$ & final speed does not matter\\
(iii) & $v^{2} = u^{2} + 2as$           & $t$ & time does not matter\\
(iv)  & $s = \dfrac{(u+v)}{2}\,t$       & $a$ & acceleration does not matter\\
\bottomrule
\end{tabular}
\end{center}
{\small\color{muted} $u$ --- initial velocity, $v$ --- final velocity, $a$ --- acceleration
(constant), $s$ --- displacement, $t$ --- time. Each problem hands a different combination
of these and asks for the rest. \textbf{Only in this copy:} the ``Missing'' column and the
$[\text{given} \to \text{find}]$ tags on the problems --- working that out himself is part
of the training.}
\else
\textbf{The toolbox} --- the four equations of uniform acceleration:
\begin{center}
\renewcommand{\arraystretch}{1.5}
\begin{tabular}{@{}ll@{}}
\toprule
(i)   & $v = u + at$\\
(ii)  & $s = ut + \tfrac{1}{2}at^{2}$\\
(iii) & $v^{2} = u^{2} + 2as$\\
(iv)  & $s = \dfrac{(u+v)}{2}\,t$\\
\bottomrule
\end{tabular}
\end{center}
{\small\color{muted} $u$ --- initial velocity, $v$ --- final velocity, $a$ --- acceleration
(constant), $s$ --- displacement, $t$ --- time. Choosing the right equation for each
problem is part of the task.}
\fi
\end{tcolorbox}

\nsec{Problems}

\begin{enumerate}[label=\textcolor{accent}{\textbf{\arabic*.}}, leftmargin=1.9em,
                  itemsep=7pt, topsep=3pt, start=0]

\item \textbf{Warm-up (done in class --- now add the graphs).}
A train enters a station at speed $v$. The station has length $L$, and the train stops
exactly at its far end, decelerating uniformly. Treating the train as a point, find the
acceleration and the stopping time. \given{given $u,\,v{=}0,\,s$ \;$\to$\; $a,\,t$}

\item \textbf{Away from the lights.}
A car starts from rest and accelerates uniformly at $2\mss$ for $6\ \mathrm{s}$.
Find its final speed and the distance covered. \given{given $u{=}0,\,a,\,t$ \;$\to$\; $v,\,s$}

\item \textbf{Between stations.}
A train speeds up from $10\ms$ to $25\ms$ in $30\ \mathrm{s}$ at constant acceleration.
Find the acceleration and the distance travelled. \given{given $u,\,v,\,t$ \;$\to$\; $a,\,s$}

\item \textbf{The cyclist.}
A cyclist moving at $6\ms$ accelerates uniformly at $2\mss$ along a $16\ \mathrm{m}$
stretch of road. Find his speed at the end of the stretch and the time it took.
\given{given $u,\,a,\,s$ \;$\to$\; $v,\,t$}

\item \textbf{Braking distance.}
A car travels at $20\ms$. The driver brakes, and the car decelerates uniformly at
$5\mss$. \;(a) How long and how far does it travel before stopping?
\;(b) If the initial speed were doubled ($40\ms$), how far would it travel? How many
times longer is that? \given{given $u,\,v{=}0,\,a$ \;$\to$\; $t,\,s$}

\item \textbf{Off the bridge.}
A stone is released from rest from a bridge and hits the water after $2\ \mathrm{s}$.
How high is the bridge, and how fast does the stone hit the water? (Ignore air
resistance.) \given{given $u{=}0,\,a{=}g,\,t$ \;$\to$\; $v,\,s$}

\item \textbf{Straight up.}
A ball is thrown straight up at $15\ms$. \;(a) How long does it take to reach the top?
\;(b) What is the maximum height? \;(c) When does it return to the hand, and with what
velocity? --- On your graphs, pay special attention: what is the acceleration \emph{at
the very top}? \given{given $u,\,a{=}-g$ \;$\to$\; $t,\,s$, symmetry}

\item \textbf{The whole trip (graphs first!).}
A bus leaves a stop: it accelerates uniformly from rest to $12\ms$ in $4\ \mathrm{s}$,
drives at a steady $12\ms$ for $6\ \mathrm{s}$, then brakes uniformly to a stop in
$3\ \mathrm{s}$. Draw all three graphs \emph{first}; then find the acceleration and the
distance for each phase, the total distance, and the average speed of the trip.
\given{three phases \;$\to$\; piecewise graphs}

\item \textbf{Read the graph.}
The $v$--$t$ graph of a tram is shown. From the graph, find the initial speed, the
acceleration (slope), and the distance travelled (area under the graph). Then sketch the
matching $a(t)$ and $x(t)$. \given{given graph \;$\to$\; $u,\,a,\,s$}

\begin{center}
\begin{tikzpicture}[x=0.95cm,y=0.42cm]
  \fill[boxbg] (0,0) -- (0,8) -- (4,0) -- cycle;
  \draw[->,thick] (0,0) -- (5.6,0) node[right=3pt]{\small $t,\ \mathrm{s}$};
  \draw[->,thick] (0,0) -- (0,9.5) node[above=3pt]{\small $v,\ \mathrm{m/s}$};
  \draw[very thick,color=accent] (0,8) -- (4,0);
  \foreach \i in {1,2,3,4,5} \draw (\i,0) -- (\i,-0.35) node[below]{\small \i};
  \foreach \j in {2,4,6,8} \draw (0,\j) -- (-0.12,\j) node[left]{\small \j};
\end{tikzpicture}
\end{center}

\item \textbf{Challenge: where equation (iv) comes from.}
Using a $v$--$t$ graph for uniform acceleration, show that the displacement equals the
area of a trapezium: $s = \tfrac{1}{2}(u+v)\,t$. Then use this to re-derive your answer
to Problem~0 in one line. \given{derivation}

\end{enumerate}

\ifdefined\isdad
% ---------------- answers (dad only) ----------------
\clearpage
\par\noindent\begin{minipage}{\linewidth}
\nsec{Answers}
{\small
\begin{enumerate}[label=\textbf{\arabic*.}, leftmargin=1.7em, itemsep=1pt, topsep=1pt, start=0]
  \item $|a| = \dfrac{v^{2}}{2L}$ directed against the motion; \; $t = \dfrac{2L}{v}$.
  \item $v = 12\ms$; \; $s = 36\ \mathrm{m}$.
  \item $a = 0.5\mss$; \; $s = 525\ \mathrm{m}$ \;(via eq.~(iv): $17.5 \times 30$).
  \item $v = \sqrt{6^{2} + 2\cdot 2\cdot 16} = 10\ms$; \; $t = 2\ \mathrm{s}$.
  \item (a) $t = 4\ \mathrm{s}$, $s = 40\ \mathrm{m}$; \;(b) $160\ \mathrm{m}$ --- $4\times$
        longer: braking distance grows as $v^{2}$ (eq.~(iii)).
  \item $h = \tfrac{1}{2}\cdot 10\cdot 2^{2} = 20\ \mathrm{m}$; \; $v = 20\ms$.
  \item (a) $1.5\ \mathrm{s}$; \;(b) $11.25\ \mathrm{m}$; \;(c) returns at $3\ \mathrm{s}$
        with $15\ms$ downward. At the top $v = 0$ but $a = -g = -10\mss$ --- the $a(t)$
        graph is one unbroken horizontal line.
  \item $a_1 = +3\mss$, $s_1 = 24\ \mathrm{m}$; \; $a_2 = 0$, $s_2 = 72\ \mathrm{m}$; \;
        $a_3 = -4\mss$, $s_3 = 18\ \mathrm{m}$. Total $114\ \mathrm{m}$ in $13\ \mathrm{s}$;
        average speed $\approx 8.8\ms$. The $v(t)$ graph is a trapezium; its area is the
        total distance.
  \item $u = 8\ms$; \; $a = -2\mss$; \; $s = \tfrac{1}{2}\cdot 8 \cdot 4 = 16\ \mathrm{m}$
        (the shaded triangle). $x(t)$ is a parabola bending down, flattening out at $t=4$ s.
  \item Trapezium with parallel sides $u$ and $v$, width $t$: area $= \tfrac{1}{2}(u+v)t$.
        For the train: $L = \tfrac{1}{2}(v+0)\,t \Rightarrow t = 2L/v$, then
        $a = (0-v)/t = -v^{2}/(2L)$. \;Same answers, no quadratic equations.
\end{enumerate}
\textcolor{muted}{Graph checklist for marking: $a(t)$ horizontal (negative where braking);
$v(t)$ straight, crossing zero exactly at the stopping/turn-around moment; $x(t)$ a
parabola with the right bend, horizontal tangent where $v=0$, never jumping.}}
\end{minipage}

% ---------------- references (dad only) ----------------
\nsec{More practice from our archive}
\begin{itemize}[leftmargin=1.2em, itemsep=3pt, topsep=2pt]
  \item \textbf{Siyavula, Physical Sciences Gr 10} ---
        {\small\path{docs/textbook/siyavula-physical-sciences/Gr10_PhysicalSciences_Learner_Eng.pdf}}\\
        {\small ch~21 \emph{Motion in one dimension}: worked examples throughout
        (PDF pp.\,403--452, graph-heavy from p.\,430) and \emph{End of chapter exercises}
        (PDF pp.\,453--457) --- includes drawing and reading $x$--$t$/$v$--$t$ graphs.
        Brief solutions in ch~25 ($\approx$ PDF p.\,504).}
  \item \textbf{NCERT, Class 9 Science} ---
        {\small\path{docs/textbook/ncert-class9-physics/ch04-describing-motion.pdf}}\\
        {\small \S4.2 graphs of motion (PDF pp.\,9--16), \S4.3 kinematic equations
        (PDF pp.\,16--18), exercises \emph{Revise, Reflect, Refine} (PDF pp.\,21--24,
        graph question on p.\,23).}
  \item \textbf{CK-12, Quizzes \& Tests} ---
        {\small\path{docs/textbook/ck12-physical-science/CK-12-Physical-Science-MS-Quizzes-and-Tests-with-Answers.pdf}}\\
        {\small quiz 12.2 \emph{Speed and Velocity} (PDF p.\,161), quiz 12.3
        \emph{Acceleration} (PDF p.\,164), Unit-4 \emph{Motion and Forces} test
        (PDF p.\,353); answers are in the same file.}
\end{itemize}
\fi

\end{document}
"
% The son's copy omits the [given -> find] tags and the missing-variable analysis:
% identifying the combination is part of the training.
\documentclass[11pt,a4paper]{article}
\usepackage[margin=2.2cm,bottom=2.4cm]{geometry}
\usepackage[T1]{fontenc}
\usepackage[utf8]{inputenc}
\usepackage{lmodern}
\usepackage{amsmath}
\usepackage{xcolor}
\usepackage{booktabs}
\usepackage{enumitem}
\usepackage{parskip}
\usepackage{fancyhdr}
\usepackage{tcolorbox}
\usepackage{tikz}
\usepackage{url}

\definecolor{accent}{HTML}{2F6FED}
\definecolor{muted}{HTML}{5B6473}
\definecolor{ink}{HTML}{1F2430}
\definecolor{boxbg}{HTML}{F2F6FF}
\definecolor{rule}{HTML}{E2E5EA}

\ifdefined\isdad
  \newcommand{\copyname}{Dad's copy: tags, answers, references}
  \newcommand{\footname}{Dad's copy}
  \newcommand{\given}[1]{\hfill{\small\color{muted}[#1]}}
\else
  \newcommand{\copyname}{student's copy}
  \newcommand{\footname}{student's copy}
  \newcommand{\given}[1]{}
\fi

\pagestyle{fancy}
\fancyhf{}
\renewcommand{\headrulewidth}{0pt}
\renewcommand{\footrulewidth}{0.4pt}
\fancyfoot[L]{\small\color{muted} Uniformly accelerated motion --- problem set 1 (\footname)}
\fancyfoot[R]{\small\color{muted} Page \thepage}

\newcommand{\nsec}[1]{\vspace{9pt}\par{\color{accent}\large\bfseries #1}\par\vspace{2pt}}
\newcommand{\ms}{\ \mathrm{m/s}}
\newcommand{\mss}{\ \mathrm{m/s^2}}

\begin{document}\color{ink}

{\huge\bfseries Uniformly Accelerated Motion}\\[4pt]
{\color{muted} Kinematics of straight-line motion with constant acceleration
\;\textperiodcentered\; problem set 1 \;\textperiodcentered\; \copyname}\\[6pt]
{\color{rule}\rule{\linewidth}{0.9pt}}

\begin{tcolorbox}[colback=boxbg, colframe=accent, boxrule=0.9pt, arc=2.5mm,
                  left=10pt, right=10pt, top=7pt, bottom=7pt,
                  title={\large\bfseries The rules of the game}, coltitle=white]
For \textbf{every} problem, including the warm-up:
\begin{enumerate}[label=\textcolor{accent}{\textbf{\arabic*.}}, leftmargin=1.8em, itemsep=1.5pt, topsep=3pt]
  \item Choose the positive direction and set $x_0 = 0$ at the start.
  \item Solve for what is asked.
  \item Draw the three graphs against time: \textbf{$a(t)$, $v(t)$, $x(t)$} --- axes labelled,
        key values marked (initial/final speeds, times, distances).
\end{enumerate}
Graphs need not be to exact scale, but the \emph{shapes} must be right: with constant
acceleration $a(t)$ is a horizontal line, $v(t)$ is a straight line with slope $a$,
and $x(t)$ is a parabola (bending up if $a>0$, down if $a<0$). Use $g = 10\mss$.
\smallskip

\ifdefined\isdad
\textbf{The toolbox} --- four equations, each \emph{missing} one quantity:
\begin{center}
\renewcommand{\arraystretch}{1.5}
\begin{tabular}{@{}llll@{}}
\toprule
& \textbf{Equation} & \textbf{Missing} & \textbf{Use when\ldots}\\
\midrule
(i)   & $v = u + at$                    & $s$ & distance does not matter\\
(ii)  & $s = ut + \tfrac{1}{2}at^{2}$   & $v$ & final speed does not matter\\
(iii) & $v^{2} = u^{2} + 2as$           & $t$ & time does not matter\\
(iv)  & $s = \dfrac{(u+v)}{2}\,t$       & $a$ & acceleration does not matter\\
\bottomrule
\end{tabular}
\end{center}
{\small\color{muted} $u$ --- initial velocity, $v$ --- final velocity, $a$ --- acceleration
(constant), $s$ --- displacement, $t$ --- time. Each problem hands a different combination
of these and asks for the rest. \textbf{Only in this copy:} the ``Missing'' column and the
$[\text{given} \to \text{find}]$ tags on the problems --- working that out himself is part
of the training.}
\else
\textbf{The toolbox} --- the four equations of uniform acceleration:
\begin{center}
\renewcommand{\arraystretch}{1.5}
\begin{tabular}{@{}ll@{}}
\toprule
(i)   & $v = u + at$\\
(ii)  & $s = ut + \tfrac{1}{2}at^{2}$\\
(iii) & $v^{2} = u^{2} + 2as$\\
(iv)  & $s = \dfrac{(u+v)}{2}\,t$\\
\bottomrule
\end{tabular}
\end{center}
{\small\color{muted} $u$ --- initial velocity, $v$ --- final velocity, $a$ --- acceleration
(constant), $s$ --- displacement, $t$ --- time. Choosing the right equation for each
problem is part of the task.}
\fi
\end{tcolorbox}

\nsec{Problems}

\begin{enumerate}[label=\textcolor{accent}{\textbf{\arabic*.}}, leftmargin=1.9em,
                  itemsep=7pt, topsep=3pt, start=0]

\item \textbf{Warm-up (done in class --- now add the graphs).}
A train enters a station at speed $v$. The station has length $L$, and the train stops
exactly at its far end, decelerating uniformly. Treating the train as a point, find the
acceleration and the stopping time. \given{given $u,\,v{=}0,\,s$ \;$\to$\; $a,\,t$}

\item \textbf{Away from the lights.}
A car starts from rest and accelerates uniformly at $2\mss$ for $6\ \mathrm{s}$.
Find its final speed and the distance covered. \given{given $u{=}0,\,a,\,t$ \;$\to$\; $v,\,s$}

\item \textbf{Between stations.}
A train speeds up from $10\ms$ to $25\ms$ in $30\ \mathrm{s}$ at constant acceleration.
Find the acceleration and the distance travelled. \given{given $u,\,v,\,t$ \;$\to$\; $a,\,s$}

\item \textbf{The cyclist.}
A cyclist moving at $6\ms$ accelerates uniformly at $2\mss$ along a $16\ \mathrm{m}$
stretch of road. Find his speed at the end of the stretch and the time it took.
\given{given $u,\,a,\,s$ \;$\to$\; $v,\,t$}

\item \textbf{Braking distance.}
A car travels at $20\ms$. The driver brakes, and the car decelerates uniformly at
$5\mss$. \;(a) How long and how far does it travel before stopping?
\;(b) If the initial speed were doubled ($40\ms$), how far would it travel? How many
times longer is that? \given{given $u,\,v{=}0,\,a$ \;$\to$\; $t,\,s$}

\item \textbf{Off the bridge.}
A stone is released from rest from a bridge and hits the water after $2\ \mathrm{s}$.
How high is the bridge, and how fast does the stone hit the water? (Ignore air
resistance.) \given{given $u{=}0,\,a{=}g,\,t$ \;$\to$\; $v,\,s$}

\item \textbf{Straight up.}
A ball is thrown straight up at $15\ms$. \;(a) How long does it take to reach the top?
\;(b) What is the maximum height? \;(c) When does it return to the hand, and with what
velocity? --- On your graphs, pay special attention: what is the acceleration \emph{at
the very top}? \given{given $u,\,a{=}-g$ \;$\to$\; $t,\,s$, symmetry}

\item \textbf{The whole trip (graphs first!).}
A bus leaves a stop: it accelerates uniformly from rest to $12\ms$ in $4\ \mathrm{s}$,
drives at a steady $12\ms$ for $6\ \mathrm{s}$, then brakes uniformly to a stop in
$3\ \mathrm{s}$. Draw all three graphs \emph{first}; then find the acceleration and the
distance for each phase, the total distance, and the average speed of the trip.
\given{three phases \;$\to$\; piecewise graphs}

\item \textbf{Read the graph.}
The $v$--$t$ graph of a tram is shown. From the graph, find the initial speed, the
acceleration (slope), and the distance travelled (area under the graph). Then sketch the
matching $a(t)$ and $x(t)$. \given{given graph \;$\to$\; $u,\,a,\,s$}

\begin{center}
\begin{tikzpicture}[x=0.95cm,y=0.42cm]
  \fill[boxbg] (0,0) -- (0,8) -- (4,0) -- cycle;
  \draw[->,thick] (0,0) -- (5.6,0) node[right=3pt]{\small $t,\ \mathrm{s}$};
  \draw[->,thick] (0,0) -- (0,9.5) node[above=3pt]{\small $v,\ \mathrm{m/s}$};
  \draw[very thick,color=accent] (0,8) -- (4,0);
  \foreach \i in {1,2,3,4,5} \draw (\i,0) -- (\i,-0.35) node[below]{\small \i};
  \foreach \j in {2,4,6,8} \draw (0,\j) -- (-0.12,\j) node[left]{\small \j};
\end{tikzpicture}
\end{center}

\item \textbf{Challenge: where equation (iv) comes from.}
Using a $v$--$t$ graph for uniform acceleration, show that the displacement equals the
area of a trapezium: $s = \tfrac{1}{2}(u+v)\,t$. Then use this to re-derive your answer
to Problem~0 in one line. \given{derivation}

\end{enumerate}

\ifdefined\isdad
% ---------------- answers (dad only) ----------------
\clearpage
\par\noindent\begin{minipage}{\linewidth}
\nsec{Answers}
{\small
\begin{enumerate}[label=\textbf{\arabic*.}, leftmargin=1.7em, itemsep=1pt, topsep=1pt, start=0]
  \item $|a| = \dfrac{v^{2}}{2L}$ directed against the motion; \; $t = \dfrac{2L}{v}$.
  \item $v = 12\ms$; \; $s = 36\ \mathrm{m}$.
  \item $a = 0.5\mss$; \; $s = 525\ \mathrm{m}$ \;(via eq.~(iv): $17.5 \times 30$).
  \item $v = \sqrt{6^{2} + 2\cdot 2\cdot 16} = 10\ms$; \; $t = 2\ \mathrm{s}$.
  \item (a) $t = 4\ \mathrm{s}$, $s = 40\ \mathrm{m}$; \;(b) $160\ \mathrm{m}$ --- $4\times$
        longer: braking distance grows as $v^{2}$ (eq.~(iii)).
  \item $h = \tfrac{1}{2}\cdot 10\cdot 2^{2} = 20\ \mathrm{m}$; \; $v = 20\ms$.
  \item (a) $1.5\ \mathrm{s}$; \;(b) $11.25\ \mathrm{m}$; \;(c) returns at $3\ \mathrm{s}$
        with $15\ms$ downward. At the top $v = 0$ but $a = -g = -10\mss$ --- the $a(t)$
        graph is one unbroken horizontal line.
  \item $a_1 = +3\mss$, $s_1 = 24\ \mathrm{m}$; \; $a_2 = 0$, $s_2 = 72\ \mathrm{m}$; \;
        $a_3 = -4\mss$, $s_3 = 18\ \mathrm{m}$. Total $114\ \mathrm{m}$ in $13\ \mathrm{s}$;
        average speed $\approx 8.8\ms$. The $v(t)$ graph is a trapezium; its area is the
        total distance.
  \item $u = 8\ms$; \; $a = -2\mss$; \; $s = \tfrac{1}{2}\cdot 8 \cdot 4 = 16\ \mathrm{m}$
        (the shaded triangle). $x(t)$ is a parabola bending down, flattening out at $t=4$ s.
  \item Trapezium with parallel sides $u$ and $v$, width $t$: area $= \tfrac{1}{2}(u+v)t$.
        For the train: $L = \tfrac{1}{2}(v+0)\,t \Rightarrow t = 2L/v$, then
        $a = (0-v)/t = -v^{2}/(2L)$. \;Same answers, no quadratic equations.
\end{enumerate}
\textcolor{muted}{Graph checklist for marking: $a(t)$ horizontal (negative where braking);
$v(t)$ straight, crossing zero exactly at the stopping/turn-around moment; $x(t)$ a
parabola with the right bend, horizontal tangent where $v=0$, never jumping.}}
\end{minipage}

% ---------------- references (dad only) ----------------
\nsec{More practice from our archive}
\begin{itemize}[leftmargin=1.2em, itemsep=3pt, topsep=2pt]
  \item \textbf{Siyavula, Physical Sciences Gr 10} ---
        {\small\path{docs/textbook/siyavula-physical-sciences/Gr10_PhysicalSciences_Learner_Eng.pdf}}\\
        {\small ch~21 \emph{Motion in one dimension}: worked examples throughout
        (PDF pp.\,403--452, graph-heavy from p.\,430) and \emph{End of chapter exercises}
        (PDF pp.\,453--457) --- includes drawing and reading $x$--$t$/$v$--$t$ graphs.
        Brief solutions in ch~25 ($\approx$ PDF p.\,504).}
  \item \textbf{NCERT, Class 9 Science} ---
        {\small\path{docs/textbook/ncert-class9-physics/ch04-describing-motion.pdf}}\\
        {\small \S4.2 graphs of motion (PDF pp.\,9--16), \S4.3 kinematic equations
        (PDF pp.\,16--18), exercises \emph{Revise, Reflect, Refine} (PDF pp.\,21--24,
        graph question on p.\,23).}
  \item \textbf{CK-12, Quizzes \& Tests} ---
        {\small\path{docs/textbook/ck12-physical-science/CK-12-Physical-Science-MS-Quizzes-and-Tests-with-Answers.pdf}}\\
        {\small quiz 12.2 \emph{Speed and Velocity} (PDF p.\,161), quiz 12.3
        \emph{Acceleration} (PDF p.\,164), Unit-4 \emph{Motion and Forces} test
        (PDF p.\,353); answers are in the same file.}
\end{itemize}
\fi

\end{document}
"
% The son's copy omits the [given -> find] tags and the missing-variable analysis:
% identifying the combination is part of the training.
\documentclass[11pt,a4paper]{article}
\usepackage[margin=2.2cm,bottom=2.4cm]{geometry}
\usepackage[T1]{fontenc}
\usepackage[utf8]{inputenc}
\usepackage{lmodern}
\usepackage{amsmath}
\usepackage{xcolor}
\usepackage{booktabs}
\usepackage{enumitem}
\usepackage{parskip}
\usepackage{fancyhdr}
\usepackage{tcolorbox}
\usepackage{tikz}
\usepackage{url}

\definecolor{accent}{HTML}{2F6FED}
\definecolor{muted}{HTML}{5B6473}
\definecolor{ink}{HTML}{1F2430}
\definecolor{boxbg}{HTML}{F2F6FF}
\definecolor{rule}{HTML}{E2E5EA}

\ifdefined\isdad
  \newcommand{\copyname}{Dad's copy: tags, answers, references}
  \newcommand{\footname}{Dad's copy}
  \newcommand{\given}[1]{\hfill{\small\color{muted}[#1]}}
\else
  \newcommand{\copyname}{student's copy}
  \newcommand{\footname}{student's copy}
  \newcommand{\given}[1]{}
\fi

\pagestyle{fancy}
\fancyhf{}
\renewcommand{\headrulewidth}{0pt}
\renewcommand{\footrulewidth}{0.4pt}
\fancyfoot[L]{\small\color{muted} Uniformly accelerated motion --- problem set 1 (\footname)}
\fancyfoot[R]{\small\color{muted} Page \thepage}

\newcommand{\nsec}[1]{\vspace{9pt}\par{\color{accent}\large\bfseries #1}\par\vspace{2pt}}
\newcommand{\ms}{\ \mathrm{m/s}}
\newcommand{\mss}{\ \mathrm{m/s^2}}

\begin{document}\color{ink}

{\huge\bfseries Uniformly Accelerated Motion}\\[4pt]
{\color{muted} Kinematics of straight-line motion with constant acceleration
\;\textperiodcentered\; problem set 1 \;\textperiodcentered\; \copyname}\\[6pt]
{\color{rule}\rule{\linewidth}{0.9pt}}

\begin{tcolorbox}[colback=boxbg, colframe=accent, boxrule=0.9pt, arc=2.5mm,
                  left=10pt, right=10pt, top=7pt, bottom=7pt,
                  title={\large\bfseries The rules of the game}, coltitle=white]
For \textbf{every} problem, including the warm-up:
\begin{enumerate}[label=\textcolor{accent}{\textbf{\arabic*.}}, leftmargin=1.8em, itemsep=1.5pt, topsep=3pt]
  \item Choose the positive direction and set $x_0 = 0$ at the start.
  \item Solve for what is asked.
  \item Draw the three graphs against time: \textbf{$a(t)$, $v(t)$, $x(t)$} --- axes labelled,
        key values marked (initial/final speeds, times, distances).
\end{enumerate}
Graphs need not be to exact scale, but the \emph{shapes} must be right: with constant
acceleration $a(t)$ is a horizontal line, $v(t)$ is a straight line with slope $a$,
and $x(t)$ is a parabola (bending up if $a>0$, down if $a<0$). Use $g = 10\mss$.
\smallskip

\ifdefined\isdad
\textbf{The toolbox} --- four equations, each \emph{missing} one quantity:
\begin{center}
\renewcommand{\arraystretch}{1.5}
\begin{tabular}{@{}llll@{}}
\toprule
& \textbf{Equation} & \textbf{Missing} & \textbf{Use when\ldots}\\
\midrule
(i)   & $v = u + at$                    & $s$ & distance does not matter\\
(ii)  & $s = ut + \tfrac{1}{2}at^{2}$   & $v$ & final speed does not matter\\
(iii) & $v^{2} = u^{2} + 2as$           & $t$ & time does not matter\\
(iv)  & $s = \dfrac{(u+v)}{2}\,t$       & $a$ & acceleration does not matter\\
\bottomrule
\end{tabular}
\end{center}
{\small\color{muted} $u$ --- initial velocity, $v$ --- final velocity, $a$ --- acceleration
(constant), $s$ --- displacement, $t$ --- time. Each problem hands a different combination
of these and asks for the rest. \textbf{Only in this copy:} the ``Missing'' column and the
$[\text{given} \to \text{find}]$ tags on the problems --- working that out himself is part
of the training.}
\else
\textbf{The toolbox} --- the four equations of uniform acceleration:
\begin{center}
\renewcommand{\arraystretch}{1.5}
\begin{tabular}{@{}ll@{}}
\toprule
(i)   & $v = u + at$\\
(ii)  & $s = ut + \tfrac{1}{2}at^{2}$\\
(iii) & $v^{2} = u^{2} + 2as$\\
(iv)  & $s = \dfrac{(u+v)}{2}\,t$\\
\bottomrule
\end{tabular}
\end{center}
{\small\color{muted} $u$ --- initial velocity, $v$ --- final velocity, $a$ --- acceleration
(constant), $s$ --- displacement, $t$ --- time. Choosing the right equation for each
problem is part of the task.}
\fi
\end{tcolorbox}

\nsec{Problems}

\begin{enumerate}[label=\textcolor{accent}{\textbf{\arabic*.}}, leftmargin=1.9em,
                  itemsep=7pt, topsep=3pt, start=0]

\item \textbf{Warm-up (done in class --- now add the graphs).}
A train enters a station at speed $v$. The station has length $L$, and the train stops
exactly at its far end, decelerating uniformly. Treating the train as a point, find the
acceleration and the stopping time. \given{given $u,\,v{=}0,\,s$ \;$\to$\; $a,\,t$}

\item \textbf{Away from the lights.}
A car starts from rest and accelerates uniformly at $2\mss$ for $6\ \mathrm{s}$.
Find its final speed and the distance covered. \given{given $u{=}0,\,a,\,t$ \;$\to$\; $v,\,s$}

\item \textbf{Between stations.}
A train speeds up from $10\ms$ to $25\ms$ in $30\ \mathrm{s}$ at constant acceleration.
Find the acceleration and the distance travelled. \given{given $u,\,v,\,t$ \;$\to$\; $a,\,s$}

\item \textbf{The cyclist.}
A cyclist moving at $6\ms$ accelerates uniformly at $2\mss$ along a $16\ \mathrm{m}$
stretch of road. Find his speed at the end of the stretch and the time it took.
\given{given $u,\,a,\,s$ \;$\to$\; $v,\,t$}

\item \textbf{Braking distance.}
A car travels at $20\ms$. The driver brakes, and the car decelerates uniformly at
$5\mss$. \;(a) How long and how far does it travel before stopping?
\;(b) If the initial speed were doubled ($40\ms$), how far would it travel? How many
times longer is that? \given{given $u,\,v{=}0,\,a$ \;$\to$\; $t,\,s$}

\item \textbf{Off the bridge.}
A stone is released from rest from a bridge and hits the water after $2\ \mathrm{s}$.
How high is the bridge, and how fast does the stone hit the water? (Ignore air
resistance.) \given{given $u{=}0,\,a{=}g,\,t$ \;$\to$\; $v,\,s$}

\item \textbf{Straight up.}
A ball is thrown straight up at $15\ms$. \;(a) How long does it take to reach the top?
\;(b) What is the maximum height? \;(c) When does it return to the hand, and with what
velocity? --- On your graphs, pay special attention: what is the acceleration \emph{at
the very top}? \given{given $u,\,a{=}-g$ \;$\to$\; $t,\,s$, symmetry}

\item \textbf{The whole trip (graphs first!).}
A bus leaves a stop: it accelerates uniformly from rest to $12\ms$ in $4\ \mathrm{s}$,
drives at a steady $12\ms$ for $6\ \mathrm{s}$, then brakes uniformly to a stop in
$3\ \mathrm{s}$. Draw all three graphs \emph{first}; then find the acceleration and the
distance for each phase, the total distance, and the average speed of the trip.
\given{three phases \;$\to$\; piecewise graphs}

\item \textbf{Read the graph.}
The $v$--$t$ graph of a tram is shown. From the graph, find the initial speed, the
acceleration (slope), and the distance travelled (area under the graph). Then sketch the
matching $a(t)$ and $x(t)$. \given{given graph \;$\to$\; $u,\,a,\,s$}

\begin{center}
\begin{tikzpicture}[x=0.95cm,y=0.42cm]
  \fill[boxbg] (0,0) -- (0,8) -- (4,0) -- cycle;
  \draw[->,thick] (0,0) -- (5.6,0) node[right=3pt]{\small $t,\ \mathrm{s}$};
  \draw[->,thick] (0,0) -- (0,9.5) node[above=3pt]{\small $v,\ \mathrm{m/s}$};
  \draw[very thick,color=accent] (0,8) -- (4,0);
  \foreach \i in {1,2,3,4,5} \draw (\i,0) -- (\i,-0.35) node[below]{\small \i};
  \foreach \j in {2,4,6,8} \draw (0,\j) -- (-0.12,\j) node[left]{\small \j};
\end{tikzpicture}
\end{center}

\item \textbf{Challenge: where equation (iv) comes from.}
Using a $v$--$t$ graph for uniform acceleration, show that the displacement equals the
area of a trapezium: $s = \tfrac{1}{2}(u+v)\,t$. Then use this to re-derive your answer
to Problem~0 in one line. \given{derivation}

\end{enumerate}

\ifdefined\isdad
% ---------------- answers (dad only) ----------------
\clearpage
\par\noindent\begin{minipage}{\linewidth}
\nsec{Answers}
{\small
\begin{enumerate}[label=\textbf{\arabic*.}, leftmargin=1.7em, itemsep=1pt, topsep=1pt, start=0]
  \item $|a| = \dfrac{v^{2}}{2L}$ directed against the motion; \; $t = \dfrac{2L}{v}$.
  \item $v = 12\ms$; \; $s = 36\ \mathrm{m}$.
  \item $a = 0.5\mss$; \; $s = 525\ \mathrm{m}$ \;(via eq.~(iv): $17.5 \times 30$).
  \item $v = \sqrt{6^{2} + 2\cdot 2\cdot 16} = 10\ms$; \; $t = 2\ \mathrm{s}$.
  \item (a) $t = 4\ \mathrm{s}$, $s = 40\ \mathrm{m}$; \;(b) $160\ \mathrm{m}$ --- $4\times$
        longer: braking distance grows as $v^{2}$ (eq.~(iii)).
  \item $h = \tfrac{1}{2}\cdot 10\cdot 2^{2} = 20\ \mathrm{m}$; \; $v = 20\ms$.
  \item (a) $1.5\ \mathrm{s}$; \;(b) $11.25\ \mathrm{m}$; \;(c) returns at $3\ \mathrm{s}$
        with $15\ms$ downward. At the top $v = 0$ but $a = -g = -10\mss$ --- the $a(t)$
        graph is one unbroken horizontal line.
  \item $a_1 = +3\mss$, $s_1 = 24\ \mathrm{m}$; \; $a_2 = 0$, $s_2 = 72\ \mathrm{m}$; \;
        $a_3 = -4\mss$, $s_3 = 18\ \mathrm{m}$. Total $114\ \mathrm{m}$ in $13\ \mathrm{s}$;
        average speed $\approx 8.8\ms$. The $v(t)$ graph is a trapezium; its area is the
        total distance.
  \item $u = 8\ms$; \; $a = -2\mss$; \; $s = \tfrac{1}{2}\cdot 8 \cdot 4 = 16\ \mathrm{m}$
        (the shaded triangle). $x(t)$ is a parabola bending down, flattening out at $t=4$ s.
  \item Trapezium with parallel sides $u$ and $v$, width $t$: area $= \tfrac{1}{2}(u+v)t$.
        For the train: $L = \tfrac{1}{2}(v+0)\,t \Rightarrow t = 2L/v$, then
        $a = (0-v)/t = -v^{2}/(2L)$. \;Same answers, no quadratic equations.
\end{enumerate}
\textcolor{muted}{Graph checklist for marking: $a(t)$ horizontal (negative where braking);
$v(t)$ straight, crossing zero exactly at the stopping/turn-around moment; $x(t)$ a
parabola with the right bend, horizontal tangent where $v=0$, never jumping.}}
\end{minipage}

% ---------------- references (dad only) ----------------
\nsec{More practice from our archive}
\begin{itemize}[leftmargin=1.2em, itemsep=3pt, topsep=2pt]
  \item \textbf{Siyavula, Physical Sciences Gr 10} ---
        {\small\path{docs/textbook/siyavula-physical-sciences/Gr10_PhysicalSciences_Learner_Eng.pdf}}\\
        {\small ch~21 \emph{Motion in one dimension}: worked examples throughout
        (PDF pp.\,403--452, graph-heavy from p.\,430) and \emph{End of chapter exercises}
        (PDF pp.\,453--457) --- includes drawing and reading $x$--$t$/$v$--$t$ graphs.
        Brief solutions in ch~25 ($\approx$ PDF p.\,504).}
  \item \textbf{NCERT, Class 9 Science} ---
        {\small\path{docs/textbook/ncert-class9-physics/ch04-describing-motion.pdf}}\\
        {\small \S4.2 graphs of motion (PDF pp.\,9--16), \S4.3 kinematic equations
        (PDF pp.\,16--18), exercises \emph{Revise, Reflect, Refine} (PDF pp.\,21--24,
        graph question on p.\,23).}
  \item \textbf{CK-12, Quizzes \& Tests} ---
        {\small\path{docs/textbook/ck12-physical-science/CK-12-Physical-Science-MS-Quizzes-and-Tests-with-Answers.pdf}}\\
        {\small quiz 12.2 \emph{Speed and Velocity} (PDF p.\,161), quiz 12.3
        \emph{Acceleration} (PDF p.\,164), Unit-4 \emph{Motion and Forces} test
        (PDF p.\,353); answers are in the same file.}
\end{itemize}
\fi

\end{document}
"
% The son's copy omits the [given -> find] tags and the missing-variable analysis:
% identifying the combination is part of the training.
\documentclass[11pt,a4paper]{article}
\usepackage[margin=2.2cm,bottom=2.4cm]{geometry}
\usepackage[T1]{fontenc}
\usepackage[utf8]{inputenc}
\usepackage{lmodern}
\usepackage{amsmath}
\usepackage{xcolor}
\usepackage{booktabs}
\usepackage{enumitem}
\usepackage{parskip}
\usepackage{fancyhdr}
\usepackage{tcolorbox}
\usepackage{tikz}
\usepackage{url}

\definecolor{accent}{HTML}{2F6FED}
\definecolor{muted}{HTML}{5B6473}
\definecolor{ink}{HTML}{1F2430}
\definecolor{boxbg}{HTML}{F2F6FF}
\definecolor{rule}{HTML}{E2E5EA}

\ifdefined\isdad
  \newcommand{\copyname}{Dad's copy: tags, answers, references}
  \newcommand{\footname}{Dad's copy}
  \newcommand{\given}[1]{\hfill{\small\color{muted}[#1]}}
\else
  \newcommand{\copyname}{student's copy}
  \newcommand{\footname}{student's copy}
  \newcommand{\given}[1]{}
\fi

\pagestyle{fancy}
\fancyhf{}
\renewcommand{\headrulewidth}{0pt}
\renewcommand{\footrulewidth}{0.4pt}
\fancyfoot[L]{\small\color{muted} Uniformly accelerated motion --- problem set 1 (\footname)}
\fancyfoot[R]{\small\color{muted} Page \thepage}

\newcommand{\nsec}[1]{\vspace{9pt}\par{\color{accent}\large\bfseries #1}\par\vspace{2pt}}
\newcommand{\ms}{\ \mathrm{m/s}}
\newcommand{\mss}{\ \mathrm{m/s^2}}

\begin{document}\color{ink}

{\huge\bfseries Uniformly Accelerated Motion}\\[4pt]
{\color{muted} Kinematics of straight-line motion with constant acceleration
\;\textperiodcentered\; problem set 1 \;\textperiodcentered\; \copyname}\\[6pt]
{\color{rule}\rule{\linewidth}{0.9pt}}

\begin{tcolorbox}[colback=boxbg, colframe=accent, boxrule=0.9pt, arc=2.5mm,
                  left=10pt, right=10pt, top=7pt, bottom=7pt,
                  title={\large\bfseries The rules of the game}, coltitle=white]
For \textbf{every} problem, including the warm-up:
\begin{enumerate}[label=\textcolor{accent}{\textbf{\arabic*.}}, leftmargin=1.8em, itemsep=1.5pt, topsep=3pt]
  \item Choose the positive direction and set $x_0 = 0$ at the start.
  \item Solve for what is asked.
  \item Draw the three graphs against time: \textbf{$a(t)$, $v(t)$, $x(t)$} --- axes labelled,
        key values marked (initial/final speeds, times, distances).
\end{enumerate}
Graphs need not be to exact scale, but the \emph{shapes} must be right: with constant
acceleration $a(t)$ is a horizontal line, $v(t)$ is a straight line with slope $a$,
and $x(t)$ is a parabola (bending up if $a>0$, down if $a<0$). Use $g = 10\mss$.
\smallskip

\ifdefined\isdad
\textbf{The toolbox} --- four equations, each \emph{missing} one quantity:
\begin{center}
\renewcommand{\arraystretch}{1.5}
\begin{tabular}{@{}llll@{}}
\toprule
& \textbf{Equation} & \textbf{Missing} & \textbf{Use when\ldots}\\
\midrule
(i)   & $v = u + at$                    & $s$ & distance does not matter\\
(ii)  & $s = ut + \tfrac{1}{2}at^{2}$   & $v$ & final speed does not matter\\
(iii) & $v^{2} = u^{2} + 2as$           & $t$ & time does not matter\\
(iv)  & $s = \dfrac{(u+v)}{2}\,t$       & $a$ & acceleration does not matter\\
\bottomrule
\end{tabular}
\end{center}
{\small\color{muted} $u$ --- initial velocity, $v$ --- final velocity, $a$ --- acceleration
(constant), $s$ --- displacement, $t$ --- time. Each problem hands a different combination
of these and asks for the rest. \textbf{Only in this copy:} the ``Missing'' column and the
$[\text{given} \to \text{find}]$ tags on the problems --- working that out himself is part
of the training.}
\else
\textbf{The toolbox} --- the four equations of uniform acceleration:
\begin{center}
\renewcommand{\arraystretch}{1.5}
\begin{tabular}{@{}ll@{}}
\toprule
(i)   & $v = u + at$\\
(ii)  & $s = ut + \tfrac{1}{2}at^{2}$\\
(iii) & $v^{2} = u^{2} + 2as$\\
(iv)  & $s = \dfrac{(u+v)}{2}\,t$\\
\bottomrule
\end{tabular}
\end{center}
{\small\color{muted} $u$ --- initial velocity, $v$ --- final velocity, $a$ --- acceleration
(constant), $s$ --- displacement, $t$ --- time. Choosing the right equation for each
problem is part of the task.}
\fi
\end{tcolorbox}

\nsec{Problems}

\begin{enumerate}[label=\textcolor{accent}{\textbf{\arabic*.}}, leftmargin=1.9em,
                  itemsep=7pt, topsep=3pt, start=0]

\item \textbf{Warm-up (done in class --- now add the graphs).}
A train enters a station at speed $v$. The station has length $L$, and the train stops
exactly at its far end, decelerating uniformly. Treating the train as a point, find the
acceleration and the stopping time. \given{given $u,\,v{=}0,\,s$ \;$\to$\; $a,\,t$}

\item \textbf{Away from the lights.}
A car starts from rest and accelerates uniformly at $2\mss$ for $6\ \mathrm{s}$.
Find its final speed and the distance covered. \given{given $u{=}0,\,a,\,t$ \;$\to$\; $v,\,s$}

\item \textbf{Between stations.}
A train speeds up from $10\ms$ to $25\ms$ in $30\ \mathrm{s}$ at constant acceleration.
Find the acceleration and the distance travelled. \given{given $u,\,v,\,t$ \;$\to$\; $a,\,s$}

\item \textbf{The cyclist.}
A cyclist moving at $6\ms$ accelerates uniformly at $2\mss$ along a $16\ \mathrm{m}$
stretch of road. Find his speed at the end of the stretch and the time it took.
\given{given $u,\,a,\,s$ \;$\to$\; $v,\,t$}

\item \textbf{Braking distance.}
A car travels at $20\ms$. The driver brakes, and the car decelerates uniformly at
$5\mss$. \;(a) How long and how far does it travel before stopping?
\;(b) If the initial speed were doubled ($40\ms$), how far would it travel? How many
times longer is that? \given{given $u,\,v{=}0,\,a$ \;$\to$\; $t,\,s$}

\item \textbf{Off the bridge.}
A stone is released from rest from a bridge and hits the water after $2\ \mathrm{s}$.
How high is the bridge, and how fast does the stone hit the water? (Ignore air
resistance.) \given{given $u{=}0,\,a{=}g,\,t$ \;$\to$\; $v,\,s$}

\item \textbf{Straight up.}
A ball is thrown straight up at $15\ms$. \;(a) How long does it take to reach the top?
\;(b) What is the maximum height? \;(c) When does it return to the hand, and with what
velocity? --- On your graphs, pay special attention: what is the acceleration \emph{at
the very top}? \given{given $u,\,a{=}-g$ \;$\to$\; $t,\,s$, symmetry}

\item \textbf{The whole trip (graphs first!).}
A bus leaves a stop: it accelerates uniformly from rest to $12\ms$ in $4\ \mathrm{s}$,
drives at a steady $12\ms$ for $6\ \mathrm{s}$, then brakes uniformly to a stop in
$3\ \mathrm{s}$. Draw all three graphs \emph{first}; then find the acceleration and the
distance for each phase, the total distance, and the average speed of the trip.
\given{three phases \;$\to$\; piecewise graphs}

\item \textbf{Read the graph.}
The $v$--$t$ graph of a tram is shown. From the graph, find the initial speed, the
acceleration (slope), and the distance travelled (area under the graph). Then sketch the
matching $a(t)$ and $x(t)$. \given{given graph \;$\to$\; $u,\,a,\,s$}

\begin{center}
\begin{tikzpicture}[x=0.95cm,y=0.42cm]
  \fill[boxbg] (0,0) -- (0,8) -- (4,0) -- cycle;
  \draw[->,thick] (0,0) -- (5.6,0) node[right=3pt]{\small $t,\ \mathrm{s}$};
  \draw[->,thick] (0,0) -- (0,9.5) node[above=3pt]{\small $v,\ \mathrm{m/s}$};
  \draw[very thick,color=accent] (0,8) -- (4,0);
  \foreach \i in {1,2,3,4,5} \draw (\i,0) -- (\i,-0.35) node[below]{\small \i};
  \foreach \j in {2,4,6,8} \draw (0,\j) -- (-0.12,\j) node[left]{\small \j};
\end{tikzpicture}
\end{center}

\item \textbf{Challenge: where equation (iv) comes from.}
Using a $v$--$t$ graph for uniform acceleration, show that the displacement equals the
area of a trapezium: $s = \tfrac{1}{2}(u+v)\,t$. Then use this to re-derive your answer
to Problem~0 in one line. \given{derivation}

\end{enumerate}

\ifdefined\isdad
% ---------------- answers (dad only) ----------------
\clearpage
\par\noindent\begin{minipage}{\linewidth}
\nsec{Answers}
{\small
\begin{enumerate}[label=\textbf{\arabic*.}, leftmargin=1.7em, itemsep=1pt, topsep=1pt, start=0]
  \item $|a| = \dfrac{v^{2}}{2L}$ directed against the motion; \; $t = \dfrac{2L}{v}$.
  \item $v = 12\ms$; \; $s = 36\ \mathrm{m}$.
  \item $a = 0.5\mss$; \; $s = 525\ \mathrm{m}$ \;(via eq.~(iv): $17.5 \times 30$).
  \item $v = \sqrt{6^{2} + 2\cdot 2\cdot 16} = 10\ms$; \; $t = 2\ \mathrm{s}$.
  \item (a) $t = 4\ \mathrm{s}$, $s = 40\ \mathrm{m}$; \;(b) $160\ \mathrm{m}$ --- $4\times$
        longer: braking distance grows as $v^{2}$ (eq.~(iii)).
  \item $h = \tfrac{1}{2}\cdot 10\cdot 2^{2} = 20\ \mathrm{m}$; \; $v = 20\ms$.
  \item (a) $1.5\ \mathrm{s}$; \;(b) $11.25\ \mathrm{m}$; \;(c) returns at $3\ \mathrm{s}$
        with $15\ms$ downward. At the top $v = 0$ but $a = -g = -10\mss$ --- the $a(t)$
        graph is one unbroken horizontal line.
  \item $a_1 = +3\mss$, $s_1 = 24\ \mathrm{m}$; \; $a_2 = 0$, $s_2 = 72\ \mathrm{m}$; \;
        $a_3 = -4\mss$, $s_3 = 18\ \mathrm{m}$. Total $114\ \mathrm{m}$ in $13\ \mathrm{s}$;
        average speed $\approx 8.8\ms$. The $v(t)$ graph is a trapezium; its area is the
        total distance.
  \item $u = 8\ms$; \; $a = -2\mss$; \; $s = \tfrac{1}{2}\cdot 8 \cdot 4 = 16\ \mathrm{m}$
        (the shaded triangle). $x(t)$ is a parabola bending down, flattening out at $t=4$ s.
  \item Trapezium with parallel sides $u$ and $v$, width $t$: area $= \tfrac{1}{2}(u+v)t$.
        For the train: $L = \tfrac{1}{2}(v+0)\,t \Rightarrow t = 2L/v$, then
        $a = (0-v)/t = -v^{2}/(2L)$. \;Same answers, no quadratic equations.
\end{enumerate}
\textcolor{muted}{Graph checklist for marking: $a(t)$ horizontal (negative where braking);
$v(t)$ straight, crossing zero exactly at the stopping/turn-around moment; $x(t)$ a
parabola with the right bend, horizontal tangent where $v=0$, never jumping.}}
\end{minipage}

% ---------------- references (dad only) ----------------
\nsec{More practice from our archive}
\begin{itemize}[leftmargin=1.2em, itemsep=3pt, topsep=2pt]
  \item \textbf{Siyavula, Physical Sciences Gr 10} ---
        {\small\path{docs/textbook/siyavula-physical-sciences/Gr10_PhysicalSciences_Learner_Eng.pdf}}\\
        {\small ch~21 \emph{Motion in one dimension}: worked examples throughout
        (PDF pp.\,403--452, graph-heavy from p.\,430) and \emph{End of chapter exercises}
        (PDF pp.\,453--457) --- includes drawing and reading $x$--$t$/$v$--$t$ graphs.
        Brief solutions in ch~25 ($\approx$ PDF p.\,504).}
  \item \textbf{NCERT, Class 9 Science} ---
        {\small\path{docs/textbook/ncert-class9-physics/ch04-describing-motion.pdf}}\\
        {\small \S4.2 graphs of motion (PDF pp.\,9--16), \S4.3 kinematic equations
        (PDF pp.\,16--18), exercises \emph{Revise, Reflect, Refine} (PDF pp.\,21--24,
        graph question on p.\,23).}
  \item \textbf{CK-12, Quizzes \& Tests} ---
        {\small\path{docs/textbook/ck12-physical-science/CK-12-Physical-Science-MS-Quizzes-and-Tests-with-Answers.pdf}}\\
        {\small quiz 12.2 \emph{Speed and Velocity} (PDF p.\,161), quiz 12.3
        \emph{Acceleration} (PDF p.\,164), Unit-4 \emph{Motion and Forces} test
        (PDF p.\,353); answers are in the same file.}
\end{itemize}
\fi

\end{document}
